\section{Video generation y agenda de investigación: formular las predicciones como preguntas verificables}

El video acumula todos los cuellos de botella anteriores: más token/latent, temporal consistency más difícil, high-resolution más costosa, y data coverage y caption quality peores. El ponente no presentó conclusiones internas sobre Sora; a partir de los fenómenos públicos, desglosó un conjunto de hipótesis de ingeniería de 2024 y, con base en ellas, propuso líneas de investigación.

\subsection{Sora-like recipe: 3D latents, scale, context y data}

La clase toma CogVideo como referencia temprana de text-to-video autoregressive y luego divide la Sora-like improvement en deflickering, image quality, resolution, long-context infrastructure y recaption/coverage. Si el latent del video tiene tamaño $T\times H\times W$, el número naive de tokens es
\[
n_v=T\frac{H}{p_h}\frac{W}{p_w},
\]
y el costo cuadrático de la full attention se dispara rápidamente; por ello, la 3D latent compression, la factorized attention, el parallelism y el data pipeline son condiciones necesarias.

\lecturefigure{slide-29-video-generation.jpg}{Video generation: de CogVideo a la Sora-like scaling recipe}{Grabación oficial de Stanford Online 00:58:29}

\begin{knowledgebox}{Qué resuelve cada uno de los cinco factores de ingeniería}
\small\sloppy
\begin{tabularx}{\textwidth}{L{0.19\textwidth} L{0.32\textwidth} X}
\toprule
Factor & Problema principal & Capacidad que no garantiza por sí solo \\
\midrule
3D latent encoder & Comprimir el espacio-tiempo y reducir el flicker & Física real y causalidad de largo plazo \\
Model/data scale & Mejorar la calidad visual y la cobertura & Ausencia de sesgo, seguridad, controlabilidad \\
High resolution & Conservar texto y detalles & Movimiento y geometría correctos \\
Long-context infra & Soportar más frames/tokens & Aprovechar efectivamente todo el context \\
Recaption y coverage & Mejorar la text-video supervision & Aprender leyes físicas a partir de raw video \\
\bottomrule
\end{tabularx}
\end{knowledgebox}

\teachervoice{00:57:54--01:04:00, el ponente califica estos factores como un resumen del Sora-like progress, no como una controlled ablation. Por eso, las notas los presentan como hypotheses verificables y no como una receta interna conocida.}

\subsection{Predicciones de 2024: reconocimiento, comprensión de video y embodied AI}

La sección anterior desglosó cómo los modelos y los datos amplían la video generation; aquí se pasa a un research forecast más incierto. La slide de predicciones sostiene que el reconocimiento visual de alto nivel se abaratará, que la video understanding será un nuevo tema candente y que la embodied AI producirá demos llamativas, pero difícilmente entrará en la vida cotidiana a corto plazo. La forma correcta de registrarlo es conservar la fecha (mayo de 2024), los supuestos de recursos y el tono de «posiblemente», en lugar de seleccionar a posteriori solo las frases que acertaron.

\lecturefigure{slide-30-research-predictions.jpg}{Predicciones de la clase en 2024: reconocimiento visual, video understanding y embodied AI}{Grabación oficial de Stanford Online 01:04:00}

\begin{warningbox}{Una prediction no es un benchmark result}
«Uno o dos años», «80\% solved» y «no podrá influir pronto en la vida real» no son criterios medibles directamente. Para evaluar una predicción hay que redefinir la task distribution, el cost, la accuracy, el long-tail, el deployment environment y el momento de la medición.
\end{warningbox}

\teachervoice{01:00:00--01:04:00, el ponente expresa a la vez optimismo y límites: el reconocimiento visual se abaratará con rapidez, pero los video models todavía tienen problemas de hallucination/counting; la embodied AI tendrá demos sorprendentes, pero podría ser costosa y difícil de desplegar en la vida diaria.}

\subsection{Recomendaciones de investigación: elegir el problema según recursos y objetivos}

La última página no es una «clasificación de temas», sino una forma de relacionar el career objective, el compute access y el research risk. Los video datasets/benchmarks son adecuados para equipos con menos recursos pero capaces de organizar la evaluation; se considera que speech tiene alta demanda pero poca inversión; architecture/optimizer puede tener gran impacto, pero exige capacidades de hardware y de sistemas; y compute-to-data es la dirección central una vez que el web corpus se va saturando.

\lecturefigure{slide-31-research-advice.jpg}{Elegir problemas de investigación según objetivos y recursos: video, speech, systems, architecture y data}{Grabación oficial de Stanford Online 01:07:34}

\begin{knowledgebox}{Reescribir las recomendaciones como matriz de decisión}
\small\sloppy
\begin{tabularx}{\textwidth}{L{0.17\textwidth} L{0.20\textwidth} L{0.20\textwidth} X}
\toprule
Dirección & Activo principal & Riesgo principal & Entregable verificable \\
\midrule
Datos/evaluación de video & collection, annotation, evaluation & Fugas, shortcut, cobertura insuficiente & dataset card, held-out test, error taxonomy \\
Speech/audio & Escenarios de usuario y datos multilingües & Privacidad, ruido, latencia en tiempo real & Métricas de ASR, TTS y de diálogo de extremo a extremo \\
Colaboración con sistemas & Topología de GPU y conocimiento de kernels & Reproducibilidad y dependencia del hardware & Throughput, latencia, utilización y costo \\
Arquitectura/optimizador & controlled runs a gran escala & Cómputo costoso, conclusiones inestables & matched-compute scaling/ablation \\
Compute-to-data & verifier, execution, search & feedback bias, synthetic collapse & provenance, verification rate, novelty \\
\bottomrule
\end{tabularx}
\end{knowledgebox}

\begin{lstlisting}[caption={Verificación previa a la publicación de un proyecto de investigación multimodal}]
freeze_task_definition_and_data_split()
record_model_version_resolution_and_prompt()
report_compute_latency_memory_and_cost()
separate_demo_success_from_repeated_trial_rate()
audit_shortcuts_leakage_and_long_tail_failures()
publish_provenance_metrics_and_negative_results()
\end{lstlisting}

\teachervoice{01:05:00--01:08:10, el ponente afirma que la speech AI está subestimada y recomienda a quienes trabajan en algoritmos colaborar cuanto antes con systems researchers, porque los mejores algoritmos deben adaptarse al hardware contemporáneo.}

\subsection{Resumen de la sección}

La video generation no consiste en ejecutar un image model durante más frames, sino en un problema de sistemas que integran la compresión espaciotemporal, el entrenamiento en paralelo, el long context, la cobertura de datos y la evaluación. Las recomendaciones de investigación también deben entenderse en función de los recursos y de los entregables verificables, en lugar de tomar las predicciones de la clase como una hoja de ruta segura.
